$\mathbf{d}_{j}$ 的元素是第 $i$ 个变量的测量值与样本均值的偏差. \hyperref[fig:3.3]{图 3.3} 显示了 $p=3$ 和 $n=3$ 时, $\mathbf{y}_{i}$ 向量分解成均值分量和偏差分量.

\begin{figure}[H]
\centering
\begin{tikzpicture}[scale=0.95]
% 三维斜投影坐标轴
\draw[->] (0,0)--(4.5,0) node[right]{$2$};
\draw[->] (0,0)--(0,3.4) node[above]{$3$};
\draw[->] (0,0)--(-1.9,-1.55) node[below left]{$1$};
\node[font=\scriptsize,below] at (1.6,0){$1$};
\node[font=\scriptsize,below] at (3.4,0){$2$};
\node[font=\scriptsize,left] at (0,1.4){$1$};
\node[font=\scriptsize,left] at (0,2.4){$2$};
\node[font=\scriptsize,above left] at (-1.0,-0.8){$1$};
% 主对角线方向（等角线）
\draw[line width=1.0pt] (0,0)--(1.5,1.15);
% y1 与 d1
\draw[->,line width=1.0pt] (0,0)--(-1.55,-1.25) node[left,font=\small]{$\mathbf{y}_{1}$};
\draw[->,line width=0.9pt] (-1.55,-1.25)--(-0.75,-0.62);
\node[left,font=\scriptsize] at (-1.15,-0.9){$\mathbf{d}_{1}$};
% y2
\draw[->,line width=1.0pt] (2.6,0.35)--(3.15,1.75) node[right,font=\small]{$\mathbf{y}_{2}$};
\node[right,font=\scriptsize] at (0.9,0.72){$\bar{x}_{1}\mathbf{1}$};
\draw[->,line width=0.9pt] (2.55,0.3)--(2.95,1.15);
\node[right,font=\scriptsize] at (2.75,0.65){$\mathbf{d}_{2}$};
% y3
\draw[->,line width=1.0pt] (1.05,0.05)--(1.65,1.35) node[right,font=\small]{$\mathbf{y}_{3}$};
\node[right,font=\scriptsize] at (1.35,-0.05){$\mathbf{d}_{3}$};
\end{tikzpicture}
\phantomsection\label{fig:3.3}\mvcap{图 3.3\quad $\mathbf{y}_{i}$ 分解为一个 $\bar{x}_{i}\mathbf{1}$ 分量和一个偏差分量 $\mathbf{d}_{i}=\mathbf{y}_{i}-\bar{x}_{i}\mathbf{1},i=1,2,3$}
\end{figure}

\medskip
\noindent\textbf{例 3.3}\quad (把一个向量分解成均值分量和偏差分量)

对例 3.2 中所给数据 $\mathbf{X}=\begin{bmatrix}4&1\\-1&3\\3&5\end{bmatrix}$, 把 $\mathbf{y}_{i}$ 分解为 $\bar{x}_{i}\mathbf{1}$ 和 $\mathbf{d}_{i}=\mathbf{y}_{i}-\bar{x}_{i}\mathbf{1},i=1,2$.

这里 $\bar{x}_{1}=(4-1+3)/3=2,\bar{x}_{2}=(1+3+5)/3=3$, 因此,
\[
\bar{x}_{1}\mathbf{1}=2\begin{bmatrix}1\\1\\1\end{bmatrix}=\begin{bmatrix}2\\2\\2\end{bmatrix}\qquad
\bar{x}_{2}\mathbf{1}=3\begin{bmatrix}1\\1\\1\end{bmatrix}=\begin{bmatrix}3\\3\\3\end{bmatrix}
\]
所以,
\[
\mathbf{d}_{1}=\mathbf{y}_{1}-\bar{x}_{1}\mathbf{1}=\begin{bmatrix}4\\-1\\3\end{bmatrix}-\begin{bmatrix}2\\2\\2\end{bmatrix}=\begin{bmatrix}2\\-3\\1\end{bmatrix}
\]
和
\[
\mathbf{d}_{2}=\mathbf{y}_{2}-\bar{x}_{2}\mathbf{1}=\begin{bmatrix}1\\3\\5\end{bmatrix}-\begin{bmatrix}3\\3\\3\end{bmatrix}=\begin{bmatrix}-2\\0\\2\end{bmatrix}
\]
注意, 因为
\[
(\bar{x}_{1}\mathbf{1})'(\mathbf{y}_{1}-\bar{x}_{1}\mathbf{1})=[2\quad2\quad2]\begin{bmatrix}2\\-3\\1\end{bmatrix}=4-6+2=0
\]
所以 $\bar{x}_{1}\mathbf{1}$ 和 $\mathbf{d}_{1}=\mathbf{y}_{1}-\bar{x}_{1}\mathbf{1}$ 是垂直的.

对于 $\bar{x}_{1}\mathbf{1}$ 和 $\mathbf{d}_{2}=\mathbf{y}_{2}-\bar{x}_{2}\mathbf{1}$, 类似的结论成立. 分解式为
\[
\mathbf{y}_{1}=\begin{bmatrix}4\\-1\\3\end{bmatrix}=\begin{bmatrix}2\\2\\2\end{bmatrix}+\begin{bmatrix}2\\-3\\1\end{bmatrix}
\]
\[
\mathbf{y}_{2}=\begin{bmatrix}1\\3\\5\end{bmatrix}=\begin{bmatrix}3\\3\\3\end{bmatrix}+\begin{bmatrix}-2\\0\\2\end{bmatrix}
\]
\hfill$\blacksquare$
